\subsection{Sprint 3 --- 06/10/2023 até 27/10/2023}

Após a retrospectiva da Sprint 2, a equipe iniciou a nova fase com uma postura mais séria e renovada. Diante dos desafios enfrentados anteriormente, encaramos a Sprint 3 como uma oportunidade de corrigir problemas de organização e avançar de forma mais consistente. Além disso, o período mais intenso de provas e trabalhos havia passado, o que permitiu maior dedicação ao projeto.

Para esta sprint, optamos por trabalhar em três \textit{user stories}. Além da continuidade da história relacionada à visualização de informações da sessão, foram incluídas tarefas de votação, cadastro e visualização de informações. Embora o backend ainda tivesse grande peso, a sprint apresentou uma distribuição mais equilibrada entre frontend e backend, o que facilitou minha participação inicial.

Na primeira metade da sprint, concentrei-me no frontend. Trabalhei na inserção de componentes em telas de gerenciamento e na criação de uma tabela genérica para essas telas. Apesar de já possuir maior familiaridade com React e Tailwind CSS, comecei a perceber que não estava totalmente satisfeito atuando apenas nessa frente. Essa percepção me levou a sair da zona de conforto e buscar envolvimento com backend.

Na parte final da sprint, assumi uma tarefa relacionada à criação de endpoints para operações de CRUD. Com o apoio de colegas mais experientes da AGES II, tive contato com conceitos importantes de desenvolvimento backend para aplicações web, como métodos HTTP, Prisma, PostgreSQL e Node.js. Até o final da sprint, consegui avançar de forma mais independente na conclusão dessa tarefa, consolidando um aprendizado técnico importante.

A equipe trabalhou intensamente durante essa sprint. Mesmo sem concluir integralmente todas as \textit{user stories}, as pendências restantes estavam relacionadas a detalhes menores. Ainda assim, voltamos a cometer problemas de gestão de tempo, deixando tarefas importantes, como o CRUD, para muito próximo do prazo final. Conseguimos contornar a situação pela dedicação da equipe, mas ficou evidente que esforço de última hora não deveria ser tratado como solução recorrente.

O encerramento da sprint ocorreu na quarta reunião com os stakeholders. A apresentação foi produtiva e atendeu às expectativas, apesar das pequenas dívidas técnicas. Esse encontro aumentou nossa confiança para a entrega final e reforçou a importância de manter comunicação clara com o cliente mesmo quando existem dificuldades internas.

Considero a Sprint 3 a mais importante para meu desenvolvimento no projeto. Ela ampliou meu conhecimento em desenvolvimento web e me permitiu descobrir uma preferência maior pelo backend. A experiência com CRUD, banco de dados e endpoints foi mais agradável do que minhas tarefas anteriores de frontend, e essa descoberta passou a influenciar minhas escolhas nas etapas seguintes e nos projetos futuros da AGES.
